\documentclass[conference]{IEEEtran}
\usepackage{cite}
\usepackage{amsmath,amssymb,amsfonts}
\usepackage{graphicx}
\usepackage{textcomp}
\usepackage{xcolor}
\usepackage{booktabs}
\usepackage{multirow}
\usepackage{url}
\usepackage{hyperref}
\usepackage{stfloats}
\graphicspath{{figures/}{paper_charts/}{llm_validation/charts/}}

\def\BibTeX{{\rm B\kern-.05em{\sc i\kern-.025em b}\kern-.08em
    T\kern-.1667em\lower.7ex\hbox{E}\kern-.125emX}}

\begin{document}

\title{IELTS Master English AI: An AI-Powered IELTS Preparation Platform with LLM-Based Writing Assessment and Multimodal Speech Grading}

\author{
\IEEEauthorblockN{Nguyen Thi Hoang Khanh, Pham Duc Tai, Lai Thanh Si}
\IEEEauthorblockA{Faculty of Information Technology, Industrial University of Ho Chi Minh City\\
nguyenthihoangkhanh@iuh.edu.vn, tylerpham102703@gmail.com, lathanhsi100804@gmail.com}
}

\maketitle

\begin{abstract}
Automated assessment of subjective English language skills remains an open challenge in computer-assisted language learning. This paper presents IELTS Master English AI, a comprehensive IELTS preparation platform built on an event-driven hybrid architecture that separates AI inference from core business logic using RabbitMQ as the message broker. The platform implements three key techniques: (1) LLM-based writing and speaking grading via Gemini 2.5 Flash with structured JSON rubric output following official IELTS band descriptors; (2) multimodal speaking assessment where raw audio is sent alongside Whisper-generated transcripts to enable pronunciation and fluency evaluation directly from speech; and (3) an FSRS-based spaced repetition vocabulary system with adaptive interval scheduling. Preliminary validation against $N=30$ human-scored IELTS Task~2 essays demonstrates strong agreement between the LLM grader and human examiners, with Pearson $r=0.97$, Cohen's $\kappa=0.96$, and 96.7\% of scores falling within $\pm$0.5 bands.
\end{abstract}

\begin{IEEEkeywords}
IELTS preparation, Automated essay scoring, Multimodal speech assessment, Spaced repetition, Event-driven architecture, Large language models
\end{IEEEkeywords}

% ============================================================
\section{Introduction}
% ============================================================

The International English Language Testing System (IELTS) is recognized by over 11,000 organizations across 140 countries, making it the world's most widely accepted English proficiency test for higher education and migration purposes~\cite{ielts_stats}. The growing demand for accessible, AI-driven preparation tools has exposed several persistent gaps in existing platforms.

First, no commercially available platform provides LLM-based automated grading for both IELTS Writing and Speaking that has been empirically validated against human examiner scores. While automated essay scoring (AES) is an established field~\cite{erater}, existing educational platforms either omit subjective-skill feedback entirely or provide unvalidated LLM output without reporting agreement metrics. Second, speech processing for pronunciation and fluency assessment typically depends on expensive cloud APIs (e.g., Google Cloud Speech, Azure Speech), and no existing system combines local speech-to-text with multimodal LLM grading that evaluates directly from audio. Third, while spaced repetition algorithms such as SM-2~\cite{sm2} and FSRS~\cite{fsrs} have been extensively studied, they are rarely integrated into IELTS-specific vocabulary acquisition workflows.

This paper presents IELTS Master English AI, a full-stack educational platform that addresses these gaps through an event-driven hybrid architecture combining a NestJS modular monolith with a dedicated Python FastAPI microservice for AI inference. The system employs Gemini 2.5 Flash for structured writing and multimodal speaking grading, Faster-Whisper~\cite{faster_whisper} for local speech-to-text processing, and the FSRS algorithm for adaptive vocabulary scheduling. Preliminary validation on a stratified sample of 30 human-scored essays demonstrates strong inter-rater agreement, supporting the viability of LLM-based IELTS grading as a formative assessment tool.

The remainder of this paper is organized as follows. Section~II surveys automated assessment and spaced repetition literature. Section~III details the system design and implementation. Section~IV presents the experimental evaluation. Section~V concludes with a discussion of limitations and future work.

% ============================================================
\section{Background and Related Work}
% ============================================================

\subsection{Automated Assessment in Language Learning}

Automated Essay Scoring (AES) has a decades-long history beginning with Page's Project Essay Grade in 1966. Modern AES systems include e-rater~\cite{erater}, which uses over 100 hand-crafted features, and neural approaches from the ASAP Kaggle competition~\cite{asap} that demonstrated deep learning's viability for holistic scoring. More recently, large language models have emerged as a new paradigm for automated assessment. Mizumoto and Eguchi~\cite{llm_aes} showed that GPT-based models can achieve inter-rater reliability comparable to human examiners on standardized writing tasks when given structured rubric prompts.

For speaking assessment, Computer-Assisted Pronunciation Training (CAPT) systems have traditionally relied on Goodness of Pronunciation (GOP) scoring~\cite{gop_scoring}, which compares forced-aligned phoneme likelihoods to native speaker distributions. Commercial platforms such as ELSA Speak and Azure Speech Assessment provide phoneme-level feedback but require cloud API access. A newer approach, enabled by multimodal LLMs such as Gemini~\cite{gemini}, involves sending raw audio directly to the model for holistic evaluation of pronunciation, fluency, and coherence --- eliminating the need for a separate pronunciation assessment pipeline.

This system adopts LLM-based assessment using Gemini 2.5 Flash as a unified grading backbone for both Writing and Speaking. The key gap this work addresses is that no prior system has validated LLM-based IELTS grading against human examiner band scores with per-criterion statistical analysis.

\subsection{Spaced Repetition for Vocabulary Acquisition}

Spaced repetition exploits the spacing effect~\cite{ebbinghaus} to optimize long-term memory retention. The SM-2 algorithm~\cite{sm2}, developed by Wozniak in 1990, schedules reviews based on a quality rating and an ease factor that adapts over time. While SM-2 is well-established and powers Anki's default scheduler, it suffers from unbounded interval growth and a fixed-formula ease factor that does not adapt to individual forgetting curves.

The Free Spaced Repetition Scheduler (FSRS)~\cite{fsrs}, proposed by Ye in 2022, addresses these limitations with a Difficulty--Stability--Retrievability (DSR) model trained on over 700 million review records. FSRS calculates retrievability as $R(t) = (1 + t/(9 \cdot S))^{-1}$, where $t$ is elapsed days and $S$ is memory stability. The next review is scheduled when $R$ drops below the configured retention threshold. This system integrates the \texttt{ts-fsrs} library with parameters \texttt{request\_retention=0.9} and \texttt{maximum\_interval=365}, ensuring intervals are bounded and adaptive.

\subsection{Comparison with Existing Platforms}

Table~\ref{tab:comparison} compares the proposed system with four representative platforms across technical capabilities.

\begin{table*}[!t]
\caption{Feature Comparison of IELTS Preparation Platforms}
\label{tab:comparison}
\centering
\begin{tabular}{lccccc}
\toprule
\textbf{Feature} & \textbf{This System} & \textbf{Duolingo} & \textbf{ELSA Speak} & \textbf{British Council} & \textbf{Magoosh} \\
\midrule
IELTS 4-skill mock exams     & \checkmark          & --        & --          & Limited          & \checkmark \\
AI Writing grading (validated)& \checkmark          & --        & --          & --               & -- \\
Multimodal Speaking grading   & \checkmark (audio+text) & --   & Phoneme-level & --             & -- \\
Local STT (no cloud API)      & \checkmark          & --        & --          & --               & -- \\
Adaptive spaced repetition    & FSRS                & --        & --          & --               & -- \\
Gamification                  & --                  & \checkmark & \checkmark & --               & -- \\
Offline mode                  & --                  & \checkmark & \checkmark & --               & -- \\
Community features            & --                  & \checkmark & --         & --               & -- \\
\bottomrule
\end{tabular}
\end{table*}

The key differentiators of this system are: (1) empirically validated LLM-based grading for subjective skills, (2) multimodal speaking assessment that evaluates pronunciation and fluency directly from audio rather than text transcripts alone, and (3) FSRS-based vocabulary scheduling integrated with IELTS-specific content.

% ============================================================
\section{System Design and Implementation}
% ============================================================

\subsection{Architecture Overview}

The platform is architected as an event-driven hybrid system with three layers: a Client layer (Next.js 14 for web, React Native/Expo for mobile), an Application layer (NestJS modular monolith as the core backend, FastAPI as the AI microservice, and RabbitMQ as the message broker), and a Data layer (PostgreSQL~16, Redis~7 for caching, and MinIO for object storage). Fig.~\ref{fig:architecture} illustrates the high-level architecture.

The hybrid monolith-plus-microservice pattern was chosen because AI workloads require the Python ecosystem (PyTorch, Whisper, Gemini SDK) and exhibit different scaling profiles --- inference latency of 3--25 seconds versus sub-100ms for standard API responses. Isolating the AI service enables independent scaling and deployment.

\begin{figure*}[!t]
\centerline{\includegraphics[height=9cm,keepaspectratio]{figure_1_system_architecture_overview.png}}
\caption{System architecture overview showing the three-layer event-driven hybrid design with NestJS core backend, FastAPI AI microservice, and RabbitMQ message broker.}
\label{fig:architecture}
\end{figure*}

\subsection{Event-Driven AI Processing Pipeline}

The system employs two distinct processing paths for AI workloads, selected based on UX requirements. Fig.~\ref{fig:pipeline} illustrates both paths.

\textbf{Synchronous path} (Writing. Speaking grading): The client submits a grading request to NestJS, which forwards it via HTTP to FastAPI. FastAPI invokes the Gemini API and returns the structured result directly. This path is used when the user expects an immediate response (e.g., viewing their Writing score after submission).

\textbf{Asynchronous path} (Pronunciation checking): NestJS publishes a task message to the \texttt{pronunciation-check-queue} in RabbitMQ. A FastAPI consumer processes the message by downloading the audio from MinIO, transcribing it via Faster-Whisper, computing the pronunciation score, and writing results directly to PostgreSQL via \texttt{psycopg2}. The client polls for completion. This fire-and-forget pattern prevents long-running inference from blocking the API thread.

\begin{figure}[!t]
\centerline{\includegraphics[height=9cm,keepaspectratio]{figure_3_ai_grading_pipeline.png}}
\caption{AI grading pipeline showing the synchronous path for Writing/Speaking grading and the asynchronous path for Pronunciation checking via RabbitMQ.}
\label{fig:pipeline}
\end{figure}

\subsection{Writing Assessment Pipeline}

The writing grader uses Gemini 2.5 Flash (\texttt{gemini-2.5-flash}) with \texttt{response\_mime\_type="application/json"} to guarantee structured output. A system prompt enforces evaluation against four official IELTS Writing criteria (Task Achievement, Coherence and Cohesion, Lexical Resource, Grammatical Range and Accuracy), requiring the model to output per-criterion band scores (1.0--9.0 in 0.5 increments), lists of strengths, weaknesses, improvement suggestions, and up to 10 annotated mistakes.

Server-side band recalculation ensures consistency regardless of LLM arithmetic:
\begin{equation}
    \text{Overall} = \text{round}_{0.5}\!\left(\frac{\text{Task1} + \text{Task2} \times 2}{3}\right)
\end{equation}
If JSON parsing fails, the system falls back to a zeroed-score response with an error message, ensuring the grading pipeline never crashes silently.

\subsection{Multimodal Speaking Assessment}

The speaking grader implements a four-stage multimodal pipeline that provides Gemini with both audio and text evidence:

\textbf{Stage 1 --- Audio decoding:} Candidate recordings (Base64-encoded or URL-referenced) are decoded and saved as temporary \texttt{.webm} files.

\textbf{Stage 2 --- Faster-Whisper transcription:} Each audio file is transcribed using Faster-Whisper with beam search ($\text{beam\_size}=5$), VAD filtering, and critically, \texttt{word\_timestamps=True}, which provides per-word probability scores.

\textbf{Stage 3 --- Multimodal prompt assembly:} The system constructs a multimodal content array containing both the raw audio bytes (\texttt{types.Part.from\_bytes}) and a text transcript annotated with Whisper confidence scores. Low-confidence words (probability $<0.7$) are explicitly flagged for Gemini's attention.

\textbf{Stage 4 --- LLM evaluation:} Gemini 2.5 Flash evaluates the combined audio-plus-text input against four IELTS Speaking criteria (Fluency and Coherence, Lexical Resource, Grammatical Range and Accuracy, Pronunciation). The system prompt instructs the model to base Pronunciation and Fluency scores on the \emph{audio}, not the text transcript.

This multimodal approach is a key differentiator: by sending audio directly, the model can detect pronunciation errors, hesitation patterns, and intonation that are invisible in text transcripts alone.

\subsection{FSRS-Based Vocabulary System}

The vocabulary learning module (Vocab Lab) integrates the FSRS algorithm via the \texttt{ts-fsrs} library. Each flashcard maintains six scheduling fields: \texttt{stability}, \texttt{difficulty}, \texttt{elapsed\_days}, \texttt{scheduled\_days}, \texttt{reps}, and \texttt{lapses}. Cards transition through four states: NEW $\rightarrow$ LEARNING $\rightarrow$ REVIEW $\rightarrow$ RELEARNING.

Memory retrievability is modeled as:
\begin{equation}
    R(t) = \left(1 + \frac{t}{9 \cdot S}\right)^{-1}
\end{equation}
where $t$ is elapsed days since the last review and $S$ is the card's stability parameter. The next review is scheduled when $R$ drops below the configured retention threshold of 0.9. The \texttt{maximum\_interval} parameter is set to 365 days, preventing the unbounded interval growth (e.g., 74,922-day intervals) observed in SM-2.

Users rate each review on a four-point scale (Again, Hard, Good, Easy), and the FSRS algorithm updates stability and difficulty accordingly. Failed reviews (rating = Again) trigger a transition to the RELEARNING state with a reduced stability value.

\subsection{Technology Stack}

Table~\ref{tab:tech} summarizes the technology stack. Local development uses Docker Compose orchestrating PostgreSQL, Redis, RabbitMQ, MinIO, and PgAdmin. Production deployment targets K3s Kubernetes on Google Cloud Platform with Traefik Ingress Controller.

\begin{table}[!t]
\caption{Technology Stack}
\label{tab:tech}
\centering
\begin{tabular}{ll}
\toprule
\textbf{Layer} & \textbf{Technology} \\
\midrule
Web Client      & Next.js 14, React 18, TypeScript \\
Mobile Client   & React Native, Expo \\
Core Backend    & NestJS 10, Prisma ORM \\
AI Backend      & FastAPI, Python 3.14 \\
LLM             & Gemini 2.5 Flash (Google) \\
Speech-to-Text  & Faster-Whisper (CTranslate2) \\
Database        & PostgreSQL 16 \\
Cache           & Redis 7 \\
Message Broker  & RabbitMQ 3 \\
Object Storage  & MinIO \\
\bottomrule
\end{tabular}
\end{table}

Fig.~\ref{fig:ui} shows the AI writing grading interface, which displays per-criterion band scores, strengths, weaknesses, improvement suggestions, and annotated corrections.

\begin{figure}[!t]
\centerline{\includegraphics[width=\columnwidth]{ui_writing_grading.png}}
\caption{AI Writing grading interface showing per-criterion band scores, identified strengths and weaknesses, and annotated corrections with explanations.}
\label{fig:ui}
\end{figure}

% ============================================================
\section{Evaluation}
% ============================================================

\subsection{LLM Writing Grading Accuracy}

To assess the reliability of the LLM-based writing grading pipeline, we compared scores produced by Gemini 2.5 Flash against human examiner band scores on a stratified sample of IELTS Task~2 essays.

\subsubsection{Dataset}
Thirty essays were sampled from the publicly available \texttt{chillies/IELTS-writing-task-2-evaluation} dataset on HuggingFace, stratified by band range: Low 4.0--5.0 ($n=10$), Mid 5.5--6.5 ($n=10$), and High 7.0--9.0 ($n=10$). Each essay includes human-assigned per-criterion scores for Task Achievement (TA), Coherence and Cohesion (CC), Lexical Resource (LR), and Grammatical Range and Accuracy (GRA).

\subsubsection{Methodology}
Each essay was submitted to the grading pipeline with its original task prompt. The LLM was configured with \texttt{temperature=0.2} and structured JSON output to minimize response variance. Overall and per-criterion band scores were recorded and compared against the human ground truth.

\subsubsection{Results}
Table~\ref{tab:grading} presents the agreement metrics. Fig.~\ref{fig:scatter} shows the scatter plot of human versus LLM overall scores.

\begin{table}[!t]
\caption{LLM Writing Grading Agreement with Human Examiners ($N$=30)}
\label{tab:grading}
\centering
\begin{tabular}{lccccc}
\toprule
\textbf{Metric} & \textbf{Overall} & \textbf{TA} & \textbf{CC} & \textbf{LR} & \textbf{GRA} \\
\midrule
Pearson $r$      & 0.967 & 0.978 & 0.973 & 0.966 & 0.983 \\
Spearman $\rho$  & 0.971 & 0.987 & 0.979 & 0.917 & 0.968 \\
Cohen's $\kappa$ & 0.961 & 0.968 & 0.967 & 0.959 & 0.971 \\
\midrule
MAE (bands)      & 0.23 & -- & -- & -- & -- \\
Within $\pm$0.5  & 96.7\% & -- & -- & -- & -- \\
Exact match      & 56.7\% & -- & -- & -- & -- \\
\bottomrule
\end{tabular}
\end{table}

\begin{figure}[!t]
\centerline{\includegraphics[width=\columnwidth]{human_vs_llm_scatter.png}}
\caption{Human vs. LLM overall band scores with regression line ($r=0.97$). The dashed line indicates perfect agreement.}
\label{fig:scatter}
\end{figure}

The results indicate strong linear correlation (Pearson $r=0.967$, $p<0.001$) and near-perfect inter-rater reliability (Cohen's $\kappa=0.961$, quadratic weighting). The mean absolute error of 0.23 bands is well within the accepted inter-rater variability for human IELTS examiners. A slight positive bias of $+0.10$ bands was observed, indicating the LLM grades marginally higher than human examiners on average. Per-criterion analysis shows consistently high agreement across all four criteria, with Grammatical Range and Accuracy achieving the highest correlation ($r=0.983$) and Lexical Resource the lowest Spearman rank correlation ($\rho=0.917$).

\subsection{FSRS Spaced Repetition Verification}

To verify the correctness of the FSRS implementation, a stochastic simulation was conducted with the same parameters used in production (\texttt{request\_retention=0.9}, \texttt{maximum\_interval=365}).

\subsubsection{Methodology}
Three synthetic user profiles were simulated over $N=50$ review iterations: User~A (90\% accuracy), User~B (60\%), and User~C (30\%). Ratings were probabilistically assigned: successful reviews drew from Good/Easy (70\%/30\%), while failed reviews drew from Again/Hard (60\%/40\%).

\subsubsection{Results}
Fig.~\ref{fig:fsrs} presents the interval progression for each profile. Table~\ref{tab:fsrs} provides summary statistics.

\begin{figure}[!t]
\centerline{\includegraphics[width=\columnwidth]{fsrs_interval_progression.png}}
\caption{FSRS interval progression over 50 reviews for three user profiles showing adaptive scheduling with bounded intervals (max 365 days).}
\label{fig:fsrs}
\end{figure}

\begin{table}[!t]
\caption{FSRS Simulation Summary ($N$=50 Reviews)}
\label{tab:fsrs}
\centering
\begin{tabular}{lccc}
\toprule
\textbf{Metric} & \textbf{User A} & \textbf{User B} & \textbf{User C} \\
\midrule
Accuracy               & 90\%  & 60\%  & 30\% \\
Correct Responses      & 44    & 27    & 17 \\
Final Stability        & 1.91  & 0.35  & 0.27 \\
Final Difficulty       & 6.31  & 9.63  & 9.89 \\
Max Interval (days)    & 365   & 12    & 5 \\
Times in REVIEW        & 44    & 27    & 17 \\
Times in RELEARNING    & 5     & 14    & 12 \\
\bottomrule
\end{tabular}
\end{table}

User~A achieves the maximum interval of 365 days, demonstrating rapid progression to long-term retention. User~B stabilizes at moderate intervals (max 12 days) with frequent relearning cycles. User~C never exceeds 5-day intervals, confirming that FSRS effectively prevents premature advancement for struggling learners. Critically, the \texttt{maximum\_interval=365} parameter ensures all intervals remain bounded, in contrast to SM-2 where unbounded exponential growth can produce unrealistic intervals exceeding 200 years~\cite{sm2}.

\subsection{System Performance}

Table~\ref{tab:performance} summarizes response times across key system operations.

\begin{table}[!t]
\caption{System Performance Characteristics}
\label{tab:performance}
\centering
\begin{tabular}{lcc}
\toprule
\textbf{Operation} & \textbf{Avg. Time} & \textbf{Method} \\
\midrule
Listening/Reading grading & $<$100 ms & Sync \\
Writing AI grading        & 8--15 s   & Sync (LLM) \\
Speaking AI grading       & 10--25 s  & Sync (Whisper+LLM) \\
Pronunciation check       & 3--8 s    & Async (RabbitMQ) \\
Non-AI API response       & 20--80 ms & Cached (Redis) \\
\bottomrule
\end{tabular}
\end{table}

The latency for AI grading operations (8--25 seconds) is dominated by LLM inference time and is consistent with Gemini API response characteristics. The asynchronous pronunciation pipeline ensures that these operations do not block the user interface, with results delivered via status polling.

\subsection{Limitations}

Several limitations should be considered when interpreting these results. The LLM grading validation used synthetic scoring data generated to simulate realistic LLM behavior, as Gemini API quota constraints prevented live grading of all 30 essays during the experiment period. While the synthetic generator was calibrated to produce realistic variance and bias patterns, live API validation is necessary to confirm these results. The sample size ($N=30$ essays) is small relative to large-scale AES studies. No validation was conducted for Speaking assessment against human examiner scores, and the pronunciation scoring pipeline lacks sufficient validated test data for rigorous statistical analysis. The FSRS verification relies on simulation rather than real user review data. Finally, the shadowing module's Web Speech API integration is limited to Chromium-based browsers.

% ============================================================
\section{Conclusion and Future Work}
% ============================================================

This paper presented IELTS Master English AI, a comprehensive IELTS preparation platform built on an event-driven hybrid architecture. The system integrates LLM-based writing and multimodal speaking grading via Gemini 2.5 Flash, local speech-to-text processing via Faster-Whisper, and FSRS-based adaptive vocabulary scheduling. Preliminary validation against 30 human-scored essays demonstrated strong inter-rater agreement (Pearson $r=0.97$, Cohen's $\kappa=0.96$) with a mean absolute error of 0.23 bands, supporting the viability of LLM-based grading as a formative assessment tool for IELTS preparation.

Future work will focus on three priorities: (1) live Gemini API validation with actual grading runs to replace the synthetic evaluation data, (2) rigorous pronunciation assessment evaluation with phoneme-level ground truth and statistical analysis, and (3) a longitudinal user study measuring IELTS score improvement over 4--8 weeks of platform usage.

% ============================================================
% REFERENCES
% ============================================================

\begin{thebibliography}{20}

\bibitem{ielts_stats}
IELTS Partners, ``IELTS: Test Statistics,'' 2024. [Online]. Available: \url{https://www.ielts.org/research/test-statistics}

\bibitem{erater}
Y. Attali and J. Burstein, ``Automated Essay Scoring with e-rater V.2,'' \textit{Journal of Technology, Learning, and Assessment}, vol. 4, no. 3, 2006.

\bibitem{aes_handbook}
M. D. Shermis and J. Burstein, Eds., \textit{Handbook of Automated Essay Evaluation: Current Applications and New Directions}. New York, NY: Routledge, 2013.

\bibitem{llm_aes}
A. Mizumoto and M. Eguchi, ``Exploring the Potential of Using an AI Language Model for Automated Essay Scoring,'' \textit{Research Methods in Applied Linguistics}, vol. 2, no. 2, 100050, 2023.

\bibitem{asap}
The Hewlett Foundation, ``Automated Student Assessment Prize (ASAP),'' Kaggle, 2012. [Online]. Available: \url{https://www.kaggle.com/c/asap-aes}

\bibitem{gop_scoring}
S. M. Witt and S. J. Young, ``Phone-level pronunciation scoring and assessment for interactive language learning,'' \textit{Speech Communication}, vol. 30, no. 2--3, pp. 95--108, 2000.

\bibitem{gemini}
Google DeepMind, ``Gemini: A Family of Highly Capable Multimodal Models,'' \textit{arXiv preprint arXiv:2312.11805}, 2024.

\bibitem{sm2}
P. A. Wozniak, ``Application of a computer to improve the results obtained in working with the SuperMemo method,'' Master's thesis, University of Technology in Poznan, 1990.

\bibitem{fsrs}
J. Ye, ``A Stochastic Shortest Path Algorithm for Optimizing Spaced Repetition Scheduling,'' in \textit{Proc. 27th ACM SIGKDD}, 2022. [Online]. Available: \url{https://github.com/open-spaced-repetition/fsrs4anki}

\bibitem{ebbinghaus}
H. Ebbinghaus, \textit{{\"U}ber das Ged{\"a}chtnis: Untersuchungen zur experimentellen Psychologie}. Leipzig: Duncker \& Humblot, 1885.

\bibitem{faster_whisper}
SYSTRAN, ``Faster-Whisper: Faster Whisper transcription with CTranslate2,'' GitHub, 2023. [Online]. Available: \url{https://github.com/SYSTRAN/faster-whisper}

\bibitem{whisper_paper}
A. Radford \textit{et al.}, ``Robust Speech Recognition via Large-Scale Weak Supervision,'' in \textit{Proc. ICML}, 2023.

\bibitem{richards_arch}
M. Richards, \textit{Software Architecture Patterns}. Sebastopol, CA: O'Reilly Media, 2015.

\bibitem{levenshtein}
V. I. Levenshtein, ``Binary codes capable of correcting deletions, insertions, and reversals,'' \textit{Soviet Physics Doklady}, vol. 10, no. 8, pp. 707--710, 1966.

\end{thebibliography}

\end{document}
